\chapter{下一阶段研究计划}
\label{chap:midterm-next-plan}

现阶段，REMUS已在双FPGA平台上完成多种处理器的运行、一致暂停、检查点保存及恢复重放。下一阶段将以这些结果为基础，完成Fabric架构下的四FPGA系统，部署双核香山昆明湖（KMH）处理器，检验已有方法在更多分区和更大规模DUT上的适用性。同时，针对前一章分析的资源失衡和执行效率问题，完善划分与运行控制方法，并通过对照实验评价改进效果。本章依次介绍平台与功能扩展、方法优化、系统评估以及论文撰写安排。

\section{四FPGA系统扩展与双核KMH部署}
\label{sec:midterm-plan-four-fpga}

\subsection{Fabric架构下的双FPGA系统验证}

已有双FPGA板级实验采用Host通过PCIe分别访问各节点的方式，而Fabric架构由Ref统一接入主机请求。因此，扩展节点数量前，首先需要在双板条件下确认这一访问方式能够承载完整的运行与状态管理流程。后续将完成Host经Ref访问本地及远端寄存器、内存的联合验证，再推进板级运行，检查控制请求与DUT跨分区数据共同经过Fabric时的执行情况。

验证将使用已有DUT和软件负载，依次覆盖程序加载、运行、定步、一致暂停、检查点保存、FPGA恢复及软件RTL重放。与原PCIe直连系统对照时，保持DUT、初始状态及输入条件一致，比较程序结果、暂停周期和检查点解码后的逻辑状态。该步骤的目标是确认控制架构变化后，系统仍保留已有功能，为四节点扩展建立可对照的基础。

\subsection{四FPGA系统与双核KMH部署}

在Fabric双板功能验证完成后，将系统扩展至四个FPGA。四个节点之间采用一跳直连的互连结构，Host仍通过Ref访问各节点的控制寄存器和内存。节点数量增加后，需要将周期协调、跨分区传播及检查点组织扩展到四个Part，使整个DUT继续作为一个逻辑仿真对象接受控制。

现有PicoRV32、Rocket、Nanhu和单核KMH均已能在双板上运行。为检验增加FPGA后对更大规模设计的承载能力，下一阶段选择双核KMH作为主要部署对象，将其划分到四个FPGA，并集成跨片执行、周期控制与扫描状态访问。软件验证将覆盖OpenSBI、Linux启动及双核负载，确认两个处理器核能够在多FPGA执行环境中协同运行。实验中将记录设计与软件版本、各Part组成及划分参数，为后续功能对照和性能分析提供一致的配置依据。

\subsection{四FPGA与复杂DUT下的功能验证}

四FPGA功能验证将围绕已有双板能力逐项展开。首先检查运行、定步与一致暂停，在不同执行位置发起暂停请求，记录各Part的进度、共同目标周期及最终停止周期。随后，将暂停点的分区状态解码并与相同逻辑周期的参考状态对照，分别验证周期对齐和状态正确性，而不只依据计数器相等判断暂停成功。

在状态管理方面，将从多个非零周期保存全局检查点，恢复四个Part及配套内存和输入，再继续执行并保存新的检查点。通过程序输出、关键执行事件和选定内部状态的比较，检查恢复后首个逻辑周期及后续区间的行为。同时，将硬件检查点加载到配套的软件RTL环境，重放相应输入并导出内部波形。上述实验将明确DUT配置、检查点周期及实际比较范围，以确认已有运行和调试功能在四FPGA与双核KMH场景下仍然成立。

\section{划分与系统性能优化}
\label{sec:midterm-plan-optimization}

\subsection{资源估计与自适应划分策略优化}

针对第\ref{sec:midterm-scalar-resource-problem}节中的资源类型失衡，后续将引入支持多维约束的划分方法，分别描述各类物理资源需求，并根据目标器件设置容量与均衡约束。资源估计将与综合、实现后的分项报告对照，分析不同模块和逻辑类型的误差，再校准估计方法。这样既能保留资源组成信息，也能检查估计值是否足以指导实际部署。

容量预算还需要考虑划分后插入的扫描链、通信和控制逻辑。后续将分别统计DUT本体与这些结构的占用，改善各Part的资源预留，减少仅按原始DUT估计而导致的容量不足。在此基础上，调整自适应层次展开策略及阈值：既避免过大的模块限制资源分配，也避免不必要的细粒度展开增加超图规模和处理时间。

\subsection{划分方法与优化效果评估}

划分实验将比较模块级、充分展开的细粒度和自适应粒度三类方案。在相同DUT、目标分区数、求解后端及资源约束下，统计超图顶点与超边数量、处理时间、峰值内存和分项资源分布，分析自适应展开能否在控制问题规模的同时取得可部署的划分结果。对无法满足容量约束的方案，也将记录超限资源及对应配置。

多维权重的效果将通过与标量方案对照评价，重点检查各类资源的最大分区负载和容量可行性，并比较校准前后的估计误差。实验将尽量在同一后端内改变权重维度；涉及KaHyPar与TritonPart之间的比较时，则分别说明后端和参数差异，避免将算法变化全部归因于资源表示。

资源均衡之外，还将记录Cut Size、最大路径切割数及对应的传播阶段数，结合实际运行开销分析划分的综合效果。最大路径切割数在此作为评价指标，用于解释划分结果对顺序传播的影响。资源分布改善是否增加通信量、较小的超图是否限制优化空间，都需要由对照结果回答，而不能仅凭某一项指标的下降判断整体性能提高。

\subsection{跨片执行与系统运行性能优化}

运行性能优化将从第\ref{sec:midterm-dual-fpga-performance-problem}节的时间分解开始。先通过双FPGA回环测量基础通路，再结合实际执行中的发送、接收、数据可用和周期推进事件，分析板间传输、片内处理及等待的耗时。进一步结合边界数据量与传播阶段数，确定限制逻辑运行频率的主要因素。

根据测量结果，将优化传输组织、缓存和流水线配置，减少关键传播链中的等待，并结合实现后的关键路径改善物理时序。对于第\ref{sec:midterm-fabric-polling-problem}节中的控制轮询问题，将推进事件通知方式的集成与验证，以减少等待状态变化时的重复读取，同时保留一致暂停所需的准确周期查询。对照时将记录控制报文数量、事件响应时间和DUT数据等待，评价通知方式对共享链路的实际影响。

各项优化均在相同DUT、划分和软件负载下比较，分别记录有效逻辑运行频率、通信等待及资源占用，并复核执行结果、一致暂停和恢复重放。物理时钟频率提高并不直接等于DUT执行速度按比例提高，因此最终评价以完成相同逻辑工作所需的实际时间为准，同时说明性能收益对应的资源代价。

\section{系统实验与对比评估}
\label{sec:midterm-plan-system-evaluation}

在完成平台功能与方法验证后，将以四FPGA REMUS为主要实验平台，使用双核KMH及已有DUT开展系统评估。前一节关注具体优化的效果，本节则评价各项机制共同运行时的执行性能、资源与状态管理开销，并通过双板和四板对照分析扩展性。

\subsection{系统执行性能评估}

执行性能实验将覆盖裸机程序、操作系统启动和应用运行等负载。每组实验记录起止逻辑周期、实际耗时及完成条件，并据此计算有效逻辑运行频率。与软件RTL仿真比较时，将保持设计、输入和工作负载一致，以完成相同执行区间所需的时间评价加速效果，并说明软件仿真器、主机配置及波形记录条件。

不同DUT的速度差异将结合划分结果、传播阶段数和跨片数据量分析。测量中分别统计持续执行时间与包含加载、暂停、检查点保存及恢复的完整流程时间，并将人工交互等待单独记录。通过这些结果，可以分别评价REMUS执行长时间负载的效率，以及实际使用检查点进行调试时的整体耗时。

\subsection{资源与状态管理开销评估}

资源评估将统计四个Part的LUT、FF、BRAM、URAM和DSP占用，利用层次化报告区分DUT、扫描与控制、通信及其他系统结构的资源需求。各组结果采用一致的统计范围，并报告实现后的时序情况。结合资源分布与关键路径，分析系统容量及物理运行频率受到的限制。

状态管理实验将测量一致暂停时间、全局检查点规模、保存时间和恢复时间。其中，扫描数据与内存镜像分别统计；暂停、扫描、主机传输和文件读写采用明确的计时边界，避免将保存命令的耗时直接视为暂停并保存的总时间。通过不同DUT和检查点配置，分析状态规模及分区分布对这些操作的影响，评价恢复重放能力带来的系统开销。

\subsection{系统可扩展性评估}

扩展性实验将使用能够在双FPGA运行的已有DUT，分别部署到双板与四板Fabric系统，在相同设计和软件负载下比较有效逻辑频率、分项资源、一致暂停时间及检查点开销。两组实验将采用一致的划分方法和硬件条件，并记录分区数量变化带来的边界数据量与最大路径切割数变化，以解释增加节点后的性能差异。

双核KMH的四FPGA部署则用于评价更大规模设计的承载能力。该实验与同一DUT的双板、四板性能对照分别回答容量扩展和运行效率问题，不以节点增加直接推定执行加速。同时，将测量双板和四板系统的控制访问延迟及状态传输开销，分析Ref统一入口和板间链路随节点数量增加的负载变化，从而说明当前架构的扩展收益及限制。

\section{论文撰写及答辩计划}
\label{sec:midterm-plan-writing}

\subsection{实验资料整理与论文撰写}

论文撰写将与系统研究同步推进。首先整理已有双FPGA实验，补齐设计与软件版本、划分参数、原始日志、波形及资源报告。四FPGA实验沿用相同的记录方式，使图表和结论能够对应到具体配置。已有记录可直接用于整理，缺少的实验条件或比较结果再安排补测。

在此基础上，完善关键的划分算法、系统实现与实验分析，说明各项改进解决的问题、实验支持的结论及适用范围。论文将区分继承自已有工作的技术、本课题的扩展与系统集成内容，并结合对照结果归纳研究贡献。形成完整初稿后，根据导师意见修订论证、相关工作和图表，补充必要的实验依据。

\subsection{阶段安排与答辩准备}

后续工作按四个阶段推进。第一阶段完成Fabric双FPGA功能验证，并整理已有板级结果。第二阶段完成四FPGA系统与双核KMH部署，验证运行、检查点和恢复重放功能。第三阶段完成划分与性能优化的对照实验，以及四FPGA系统的性能、开销和扩展性评估。第四阶段集中完成论文修订，准备预答辩和正式答辩材料。

实验整理与正文写作贯穿各阶段，随研究进展同步更新。答辩材料围绕研究问题、关键方法和实验证据展开，并核对论文、图表及演示中的配置、数据和完成状态，保持表述一致。
